%% kapa: exam=1 round=29 year=2018 period=2 form=A subject=accounting subject_ko=회계학 seq=10 number=50
%% source: Q-Net 감정평가사 기출문제 (cst00309) artlSeq=5202808 file=extracted/1008/5202808/2018년 제29회 감정평가사 1차시험 기출문제/2018년 제29회 감정평가사 1차시험 기출문제/2교시 A형.pdf
%% conversion: kordoc 4.24.51 -> tools/gen1.py (auto); review=proofread-sonnet; flags=none
\begin{question}{50}
다음 설명 중 옳은 것을 모두 고른 것은?
\begin{framedtext}
ㄱ. 특정 유형자산을 재평가할 때, 해당 자산이 포함되는 유형자산 분류 전체를 재평가한다. \par
ㄴ. 자가사용부동산을 공정가치로 평가하는 투자부동산으로 대체하는 시점까지 그 부동산을 감가상각하고, 발생한 손상차손을 인식한다. \par
ㄷ. 무형자산으로 인식되기 위해서는 식별가능성, 자원에 대한 통제 및 미래경제적효익의 존재 중 최소 하나 이상의 조건을 충족하여야 한다. \par
ㄹ. 무형자산을 창출하기 위한 내부 프로젝트를 연구단계와 개발단계로 구분할 수 없는 경우에는 그 프로젝트에서 발생한 지출은 모두 개발단계에서 발생한 것으로 본다. \par
\end{framedtext}
\begin{choices}
  \item ㄱ, ㄴ
  \item ㄱ, ㄷ
  \item ㄴ, ㄹ
  \item ㄷ, ㄹ
  \item ㄱ, ㄴ, ㄷ
\end{choices}
\end{question}
\answer{1}
